\section{Coned geometry of graphical presentations}\label{sec:geometry}

This section establishes the geometric control needed by the module and
spectral arguments: relative to one root, all intersections of a translated
relator with relators no farther from the root lie in one short intrinsic
ball. We first define the labelled graphs and their coned Cayley graph,
then prove an avoidance estimate and uniform thinness. The existence of
the graph data, with the required expansion properties, is established in
Section~\ref{sec:data}.

Let \((\Theta_i)\) be finite connected simple \(d\)-regular graphs,
\(d\ge3\), of girths \(\sigma_i\). Label their oriented edges by a finite
alphabet \(S^{\pm1}\), with reversal corresponding to formal inversion.
Assume that the labelling is reduced, meaning locally injective into the
bouquet with alphabet \(S\), and that
\begin{equation}\label{eq:geometric-parameters}
0<\lambda\le\frac1{12},\qquad \lambda\sigma_i>3.
\end{equation}
We require the following uniqueness property: the label word of a simple
path in \(\Theta_i\) of length at least \(\lambda\sigma_i\) occurs on no
other simple path in the disjoint family.

Define
\[
\Gamma=\langle S\mid\text{labels of all closed walks in all }\Theta_i\rangle,
\qquad X=\Cay(\Gamma,S).
\]
Initially \(X\) is the formal labelled Cayley graph: it has the specified
edge from \(h\) to \(hs\) for each \(s\in S\), and the reverse traversal
has label \(s^{-1}\). This graph is deterministic for reading words, even
before loops or multiple edges have been excluded. Assume that the based
label maps \(\phi_i:\Theta_i\to X\) are isometric embeddings on vertices.
They are well-defined by the presentation. For every \(g\in\Gamma\), write
\[
P_{g,i}=g\phi_i(\Theta_i)
\]
for the corresponding \emph{copy}, with its source identification, and put
\(\sigma(P_{g,i})=\sigma_i\). Denote intrinsic distance in a copy by
\(d_P\).

\begin{lemma}\label{lem:distinct-copies}
Different indexed copies cannot share a simple edge path of length at
least \(\lambda\sigma(P)\) in either copy \(P\). In particular, different
index pairs give different labelled subgraphs.
\end{lemma}
\begin{proof}
Read a common path in the same direction in \(X\). Its pullbacks are
simple, since the vertex maps are injective. Label uniqueness identifies
the source paths and hence their source graph. Equality at the first
target vertex then identifies the translating elements. The graphs have
simple paths longer than the threshold, by regularity and girth, so
equality of whole copies also forces equality of their indices.
\end{proof}

Adjoin a new center \(c_P\) for every copy and join it by an edge to every
vertex of \(P\). The resulting graph is denoted \(\widehat X\); all edges
have length one.

\subsection{Diagrams with adjustable sides}

We shall use one boundary estimate twice: first to control paths avoiding
a specified center, and then to obtain short fillings in \(\widehat X\).
For a passage through a center, we hold fixed its endpoints and assigned
copy, while allowing the path inside that copy to change. A minimal
filling will cut these paths into linearly many arcs, counted in terms
of the length of the coned walk. An arc separating different assigned
copies will be short relative to the girth of either copy.

Start a nonempty closed edge walk in \(\widehat X\) at a group vertex
and partition it into \(m\) sides. Each side is a generator edge in
\(X\) or a two-edge passage through a center. Inflate a center passage
to any walk between its endpoints in its assigned copy. Such a side is
\emph{adjustable}; an empty walk is allowed when the endpoints coincide.
Generator sides remain fixed. The result is a closed walk in \(X\).

For an inflated loop, use a finite connected plane multigraph diagram
whose exterior contour maps to that loop and whose bounded-face contours
are closed walks in assigned copies. Such diagrams exist by van Kampen's
lemma: the presentation allows all closed walks in the source graphs,
and the group value at a face's initial vertex selects its translated
copy. Contours may repeat vertices and edges. This is essential for the
surgeries below: deleting a common edge merges its incident regions,
but other contacts can remain as repeated vertices or bridges. The
merged bounded contour must remain a closed walk in its assigned copy;
an adjusted exterior side must remain a walk with its prescribed
endpoints and assignment.

An \emph{incidence} is one occurrence of an edge on a contour; a bridge
has two incidences on the same contour. The exterior contour retains its
partition into side occurrences. Mark the positions where sides end,
and also the corresponding plane-graph vertices. There are at most
\(m\) distinct marked vertices, although the contour may visit a marked
vertex repeatedly. Only the marked positions are side breaks: another
visit to the same vertex need not end the current side. If inflation
gives the empty walk, an isolated marked vertex is permitted.

Minimize the number \(F\) of bounded faces, then the number of edges,
over diagrams and adjustable walks with the prescribed side data. Suppress
all unmarked degree-two vertices; the resulting edges will be called
\emph{arcs}, and their number will be denoted \(b\). An arc retains its
original edge length. Loop arcs are allowed.

\begin{lemma}[Short arcs and their number]\label{lem:diagram-count}
In a diagram chosen as above, the following hold.
\begin{enumerate}[label=\textup{(\alph*)}]
\item Every bounded face assigned to \(P\) has perimeter at least
\(\sigma(P)\). Every arc incident to that face has original length less
than \(\lambda\sigma(P)\).
\item More generally, an arc with different copy assignments on its two
incidences has length less than \(\lambda\sigma(Q)\) for either
assignment \(Q\). An arc with a generator-side incidence has length one.
\item The compressed diagram satisfies
\begin{equation}\label{eq:diagram-count}
(1/\lambda-6)F\le6m-6,\qquad b\le6m.
\end{equation}
\end{enumerate}
\end{lemma}

\begin{proof}
We first record the consequences of minimality, with contour incidences
understood throughout.

\emph{Face words and mergers.}
A bounded-face word cannot be freely trivial. Indeed, the diagram
identity expresses the exterior word as a product of conjugates of the
face words, each used once. Omitting a freely trivial factor and applying
van Kampen's lemma again gives a diagram with fewer faces for the same
exterior word. Its prescribed exterior paths lift identically by
determinism of \(X\). A closed walk in \(P\) shorter than \(\sigma(P)\)
cancels to the empty walk, so every face has the claimed lower perimeter
bound.

Suppose two distinct bounded faces with the same copy assignment share
an edge. That edge is not a bridge. Delete it and join the two contours
along its former incidences. Every remaining step lies in the same copy,
so the merged contour is an admissible closed walk there. Exactly one
bounded face disappears. Other contacts between the old contours may
produce repeated vertices or bridges, which our diagram category allows.
Thus this configuration contradicts minimality of \(F\).

The same operation excludes an edge shared by a bounded face and an
adjustable exterior side with the same assignment. Replace this
particular occurrence of the edge in the side by the rest of the
bounded contour. The resulting side is still a walk in that copy with
the same endpoints. Its prescribed position among the exterior sides
is unchanged, and again one bounded face disappears.

\emph{Leaves and bounded-face bridges.}
There is no unmarked vertex of valence one. Its leaf edge gives an
immediate backtrack; on the exterior this lies within one adjustable side,
because a side break would be marked and a generator side has marked
endpoints. Erasing the leaf preserves the assignments and lowers the edge
count.

There is also no bridge whose two incidences belong to a bounded face.
After its open edge is removed, the exterior contour lies wholly in one
component. Keep that component and discard the other component together
with the bridge. On the affected bounded contour this replaces
\(q e\gamma e^{-1}q'\) by \(q q'\), where \(e\) is the bridge and
\(\gamma\) is the entire excursion around the discarded component.
The retained walk is still closed and still lies in the assigned copy;
no free triviality of \(\gamma\) is required. The exterior, including
its side data, is unchanged. The face count cannot increase---also
directly from \(F=E-V+1\)---and the edge count decreases.

\emph{Short arcs.}
The incidence type on either side of a compressed arc is constant: it is
a bounded face or an exterior side of the prescribed decomposition.
Changes of exterior side occur at marked vertices and are not suppressed.
An arc with a generator-side incidence therefore has length one. By the
preceding restrictions, an arc incident to a bounded face has either such
an incidence on the other side or a different copy assignment there.

Consider an arc with different assignments and length greater than two.
Its word is freely reduced. Otherwise a consecutive cancelling pair has
an unmarked bivalent midpoint and two distinct endpoints in the diagram.
The latter assertion holds even for a loop arc: its internal vertices are
distinct and closure occurs only at the two ends of the entire arc.
Contract the cancelling pair as a length-two tree, identifying its outer
endpoints, which have the same group value. At each incident contour the
two steps cancel, so assignments and side data are preserved. This
non-loop contraction preserves \(F\) and lowers the edge count, a
contradiction.

If the arc had length at least \(\lambda\sigma(Q)\) for one assignment
\(Q\), its first \(\lceil\lambda\sigma(Q)\rceil\) edges would be a
nonbacktracking path shorter than the girth. It is simple in \(Q\), hence
also in the other source by the vertex embeddings. This contradicts
Lemma~\ref{lem:distinct-copies}. Arcs of length at most two are already
shorter than the threshold. This proves (a) and (b).

\emph{Counting.}
There is a marked vertex, so suppression cannot leave an exceptional
unmarked circle component. Every unmarked surviving vertex has valence
at least three. If \(v_D\) is the compressed vertex count, then
\[
v_D-b+F=1,\qquad 2b\ge3(v_D-m),
\qquad b\le3F+3m-3.
\]
By (a), every bounded face has at least \(1/\lambda\) arc incidences,
counted with multiplicity. Their total is at most \(2b\), so
\[
F/\lambda\le2b\le6F+6m-6.
\]
This is the first inequality in \eqref{eq:diagram-count}. Since
\(\lambda\le1/12\), it gives \(F\le m-1\) and then \(b\le6m\).
The same formulas cover trees and the isolated marked vertex.
\end{proof}

\subsection{Avoidance and uniform thinness}

The arc bound has two geometric consequences. Closing an avoiding path
by a designated cone passage bounds its endpoints' intrinsic distance.
Replacing arcs by cone passages while retaining fixed generator edges
produces a linear filling inequality with universal constants.

\begin{proposition}\label{prop:cone-geometry}
There is an integer \(D_0\), independent of all the above data, such that
every vertex on a geodesic triangle in \(\widehat X\) is within \(D_0\) of
a vertex on another side. Moreover:
\begin{enumerate}[label=\textup{(\roman*)}]
\item\label{item:avoidance}
If an edge path of length \(L\) joins \(x,a\in P\) while avoiding
\(c_P\), then
\begin{equation}\label{eq:avoidance}
d_P(x,a)\le12(L+1)\lambda\sigma(P).
\end{equation}
\item\label{item:simple-cayley}
If \(\lambda<1/12\), the graph \(X\) is simple: no nonfreely-trivial
closed label word in \(X\) has length at most two.
\end{enumerate}
\end{proposition}

\begin{proof}[Proof of Proposition~\ref{prop:cone-geometry}]
\emph{Avoidance.}
For \eqref{eq:avoidance}, close the given avoiding path by one designated
side through \(c_P\). Then \(m\le L+1\), and no other side has center
\(c_P\). Use a minimal diagram. An arc incidence on the designated side
is opposite a bounded face of different assignment or an exterior
incidence across a bridge. In the latter case the other incidence cannot
also lie in the designated side. If it did, the interval from that
bridge's first traversal to its retraversal in the \emph{linear} side
walk would contain no prescribed side break. Excise the visited component
and the bridge. This shortens only the designated side, preserving its
endpoints and copy assignment, and cannot increase \(F\). It contradicts
minimality. This argument concerns marked positions; additional visits to
a marked vertex do not insert a side break.

Thus every designated-side incidence has length less than
\(\lambda\sigma(P)\), either by Lemma~\ref{lem:diagram-count}(b) or
because its opposite incidence is a generator side. There are at most
\(2b\) such occurrences. The adjusted walk inside \(P\) has length at
most \(12(L+1)\lambda\sigma(P)\), giving the distance bound. Empty
walks, including the case \(x=a\), satisfy it as well.

\emph{Simplicity of the Cayley graph.}
For part~\ref{item:simple-cayley}, consider a closed word of one or two
generator sides. Strict \(\lambda<1/12\) in
\eqref{eq:diagram-count} forces \(F=0\). The exterior contour of a tree
is freely trivial. This excludes identity generators, unintended
coincidences of generator edges, and generator involutions, as claimed.

\emph{Short fillings.}
To establish uniform thinness, we first show that
every edge loop of length \(n>0\) in \(\widehat X\) has a filling with
at most \(12n\) relation faces of perimeter at most four. For its
\(m\le n\) sides, take the compressed minimal diagram. Choose one
\emph{standard path} for each arc, consistently under reversal. Keep the
unit edge if the arc has a generator-side incidence. Otherwise an
assignment at an incidence contains both endpoints, so use the two-edge
path through that copy's center.

Around a bounded face assigned to \(P\), replace each standard path by
the two-edge path through \(c_P\). Each replacement uses one closed walk
of length at most four. The resulting contour freely contracts through
that one center. Thus the standard exterior walk can be filled at cost
at most the number of bounded-face incidences. On an exterior adjustable
side, make the same replacements through its prescribed center. The
consecutive pairs reduce to the required single cone passage. If an
inflated side was empty, insert its desired backtrack freely. Generator
sides need no change. The total cost is at most the number of bounded-face
and exterior adjustable-side incidences, hence at most
\(2b\le12m\le12n\).

This filling argument works for the singular diagrams being used. More
explicitly, for directed edge paths with their endpoint values, the
plane-diagram identity expresses the exterior contour as a product of
conjugates of the bounded contours, one for each face, up to free edge
cancellations and their inverses. Substituting the indicated fillings
gives the claimed product of short relations. Conversely, the van Kampen
construction from this product gives a plane diagram over the graph:
inverse-edge folds identify only vertices with the same target value.

\emph{From fillings to thin triangles.}
The bound implies the coarse filling hypothesis in
\cite[III.H.2.1 and III.H.2.9]{BH1999} with universal constants. Subdivide a rectifiable
loop of length \(\ell\) into \(O(\ell+1)\) arcs of length at most one.
Choose nearby vertices at distance at most one and join successive ones
by paths of at most three edges. Fill this edge loop as above. In a
resulting small-face diagram with \(f\) faces and exterior length \(n\),
\[
2E=n+\sum_F|\partial F|\le n+4f,\qquad V\le E+1.
\]
Thicken its finite plane graph to a regular neighborhood, using vertex
disks and edge strips, and cap the bounded contour components. The result
is a disk. The vertex regions are mapped to vertices, the strips along
their edges, and each cap has a boundary of bounded image perimeter.
Triangulation has total cost \(O(E+V+f)\); each triangle can be assigned a
map of uniformly bounded image diameter. Coarse fillings do not require
continuous maps on their interiors. An additional collar with \(O(n)\)
triangles restores the prescribed exterior parameterization: each block
joins a unit-edge traversal to that traversal with its inserted vertex
pauses, and its image lies in the same edge. The short paths to the original
subdivision points give an annulus with \(O(\ell+1)\) further triangles
of bounded image diameter. Constant loops or empty approximating walks
are filled separately by a bounded number of triangles.

The constants in this construction are independent of the graphs,
alphabet, and \(\lambda\). A connected unit-edge graph is geodesic,
regardless of local finiteness. The cited theorem therefore gives a
hyperbolicity constant depending only on these universal coarse filling
constants. Enlarging it to an integer accounts for moving points on
triangle sides to vertices and gives the required \(D_0\).
\end{proof}

\subsection{Simultaneous overlaps}

Set, once and for all,
\begin{equation}\label{eq:root-constants}
t=D_0+3,\qquad L_0=2+2t+2D_0.
\end{equation}
The next consequence is stronger than a pairwise intersection bound: one
ball controls all the intersections in the specified distance order.

\begin{lemma}[Simultaneous overlap bound]\label{lem:overlap}
Let \(R\) be a group vertex or a center of \(\widehat X\), and let
\(A=c_P\ne R\). Choose one geodesic from \(A\) to \(R\), whose first
group vertex is \(a\in P\). Every vertex shared by \(P\) and any
\emph{different} copy \(Q\) with
\[
d_{\widehat X}(c_Q,R)\le d_{\widehat X}(c_P,R)
\]
lies in the intrinsic induced ball \(B_P\) about \(a\) of radius
\begin{equation}\label{eq:overlap-radius}
r_P=\left\lfloor12(L_0+1)\lambda\sigma(P)\right\rfloor.
\end{equation}
\end{lemma}

\begin{proof}
Take such a \(Q\) and a shared vertex \(x\), and write \(B=c_Q\).
The path \(A,x,B\) is a length-two geodesic. A geodesic from \(B\) to
\(R\) avoids \(A\), since passing through \(A\) would give length
\(2+d(A,R)>d(A,R)\). All distances in the rest of this paragraph are in
\(\widehat X\).

If \(d(A,R)\ge t\), let \(z\) be the vertex at distance \(t\) along the
chosen \(A\)-to-\(R\) geodesic. Its distance from every vertex of
\(A,x,B\) is at least \(t-2>D_0\). Thinness supplies a vertex \(y\) on
the \(B\)-to-\(R\) side with \(d(y,z)\le D_0\). A shortest bridge
between them avoids \(A\), because \(d(A,z)=t>D_0\). Follow
\(x,B,y,z,a\), using that bridge and the indicated geodesic segments.
This walk avoids \(A\), and its length is at most
\[
1+(2+t+D_0)+D_0+(t-1)=L_0.
\]
If \(d(A,R)<t\), use instead the walk from \(x\) through \(B\) to
\(R\) and back to \(a\). It avoids \(A\) and has length at most
\(1+d(B,R)+d(A,R)-1\le2d(A,R)\le L_0\).

Apply \eqref{eq:avoidance} to obtain
\(d_P(x,a)\le12(L_0+1)\lambda\sigma(P)\). Intrinsic vertex distances
are integers, so taking the floor gives the claimed radius. Crucially,
the chosen \(a\) does not depend on \(Q\) or \(x\).
\end{proof}
